\section{Constraint ranks and route controls}

A constant symmetric $n$ can be diagonalized by $SO(3)$. The $A_k$ in (H-03) represent the adjoint structure algebra; the Jacobi identity yields the momentum Poisson algebra. Both $T$ and $P$ are invariant under this action, giving $\{J,F\}=0$. In the six coordinates and their conjugates these statements are directly checkable polynomial identities.

\stepcounter{equation}\begin{equation}\tag{CC-01}\label{eq:CC-01}
\begin{aligned}g&=\mathrm{diag}(1,2,3),\quad c=1,\quad M=\pi g,\\ M_{T^3}&=\mathrm{diag}(1,1,0),\quad M_{\rm Nil}=\sqrt3 I,\\ M_{\rm Sol}&=\frac{5\sqrt3}{3}I,\quad M_{S^3}=\mathrm{diag}(4\sqrt2,0,0).\end{aligned}
\end{equation}

All representatives satisfy $F=J=0$. For $T^3$ the minor is $\partial F/\partial p_{33}=-6$. For the others take the required initial rows of $(F,J_1,J_2,J_3)$ and columns of $(p_{11},p_{23},p_{13},p_{12})$. Their determinants are $9\sqrt3$, $50\sqrt3$, and $-64\sqrt2$ for Nil, Sol and $S^3$, respectively; the Sol value uses the diagonalized $n=\mathrm{diag}(0,1,-1)$. Thus the constraint Jacobian ranks are $1/3/4/4$.

The spans of the $A_k$ have dimensions $0/2/3/3$, bounding the ranks of $J$. Near the nonzero minors, these maximal ranks and the independence of $dF$ persist, defining the regular strata of (H-04). Counts are not extended uniformly to points where minors vanish, such as symmetry-enhanced configurations.

\subsection{Closure for quadratic actions}

\stepcounter{equation}\begin{equation}\tag{CC-02}\label{eq:CC-02}
M_{A b}=\{\Phi_A,\phi_b\},\quad u^AM_{Ab}=0\quad\Longrightarrow\quad u^A\{\Phi_A,H_c\}\simeq0.
\end{equation}

Here $\phi$ denotes the original primaries and $\Phi$ all independent constraints available at the current stage. Add only new independent linear constraints from the left nullspace of $M$, iterating until none remain. For constant-coefficient quadratic actions this linear-algebra procedure is exact. It is not used to replace nonlinear constraints by linear-coefficient ranks. If the final Poisson rank is $s$ and there are $m$ independent constraints, their first- and second-class counts are $m-s$ and $s$.

\begin{center}\small\begin{tabularx}{\textwidth}{@{}>{\raggedright\arraybackslash}X>{\raggedright\arraybackslash}X>{\raggedright\arraybackslash}X@{}}\toprule
Control & Connection: FC/SC/pairs & Eliminated: FC/SC/pairs\\\midrule
Complete torus quadratic & 1 / 4 / 1 & 1 / 0 / 1\\
Complete sphere quadratic & 0 / 8 / 0 & 0 / 4 / 0\\
\bottomrule\end{tabularx}\end{center}

For the eliminated sphere block, $H=-p_u^2/24-6\ell u$ and $p_\ell=0$ give $u=0$, then $p_u=0$, then $\ell=0$. On the torus the lapse primary is first class and the quadratic action retains a scalar pair. Section 5 treats the condition from the original nonlinear constraint.


\subsection{Scalar controls from the complete action}

Set $c=1$, $n=2kI$ and write the independent scalar connection as $\omega^0{}_i=b e^i$, $\omega^j{}_k=a e^i$ cyclically. Here $k=0$ is the torus and $k=1$ the round-sphere reference.

\stepcounter{equation}\begin{equation}\tag{D-02}\label{eq:D-02}
L_{\rm P}=6Nq^3b^2-12q^2\dot q\,b+12kNq^2a-6Nq^3a^2,\quad a=k/q,\quad b=\dot q/(Nq).
\end{equation}

\stepcounter{equation}\begin{equation}\tag{D-03}\label{eq:D-03}
\begin{aligned}L_{\rm met}&=6k^2Nq-6q\dot q^2/N,\\ q&=1+\varepsilon u,\quad N=1+\varepsilon\ell,\quad a=k+\varepsilon A,\quad b=\varepsilon B,\\ L^{[1]}&=6k^2(\ell+u),\\ L_{\rm P}^{[2]}&=-6A^2-12kAu+6B^2-12B\dot u+6k^2\ell u-6k^2u^2,\\ L_{\rm met}^{[2]}&=-6\dot u^2+6k^2\ell u.\end{aligned}
\end{equation}

Substitution of the quadratic connection equations $A=-ku$, $B=\dot u$ gives the same metric action. The sphere has a nonzero first variation; the zero-pair agreement is an off-shell algebraic control. The torus quadratic scalar action retains one pair, while the original lapse equation requires $\dot u^2=0$, the pure-scale section of the cone in Section 5.


\subsection{Counts at symmetric Taub points and nonzero modes}

At the Taub point (BF-01), the homogeneous momentum and total ranks are two and three, leaving three pairs in the linear quotient. This is a different stratum from the regular representative (CC-01), with ranks three and four. Appendix G gives the $j=0,1$ KID exceptions and nonzero minors valid on the whole interval. The 40/34/10-configuration flat comparison in Section 4 counts one real Fourier quadrature, separately from the regular homogeneous count (H-04).
